\documentclass[a4paper,10pt]{article}

\usepackage{url}

\usepackage[utf8]{inputenc}
\usepackage[T1]{fontenc}
\usepackage[spanish, es-tabla]{babel}
\usepackage[indent]{parskip}

\clubpenalty=10000
\widowpenalty=10000

\usepackage[gen]{eurosym}

\usepackage{hyperref}
\hypersetup{
    colorlinks=true, % make the links colored
    linkcolor=blue, % color TOC links in blue
    urlcolor=black, % color URLs in red
    linktoc=all % 'all' will create links for everything in the TOC
}

\usepackage[dvipsnames]{xcolor}
\usepackage{setspace}
\usepackage{float}


\usepackage{biblatex}
\bibliography{fuentes.bib}


\setlength{\parindent}{0cm}
\usepackage{amsmath, amssymb}

\usepackage{csquotes}

\renewcommand{\it}[1]{\textit{#1}}
\renewcommand{\bf}[1]{\textbf{#1}}
\renewcommand{\tt}[1]{\texttt{#1}}

\newcommand{\car}[1]{`\texttt{#1}'}
\newcommand{\tok}[1]{\texttt{\textless#1\textgreater}}





% Listados de código
\usepackage{listings}
\lstdefinelanguage{minileng}{
  %keyword1&2&6
  morekeywords = [1]{programa, principio, fin, si, ent, si\_no, fsi, mq, fmq, accion},
  %keyword3
  morekeywords = [2]{escribir, leer, entacar, caraent},
  %keyword4
  morekeywords = [3]{entero, caracter, booleano},
  %keyword5
  morekeywords = [4]{ref, val},
  %function1
  sensitive = true,
  morecomment = [l]{\%},
  morecomment = [s]{\%\%}{\%\%},
  % morecomment = [s]{/**}{*/},
  morestring = [b]",
  stringstyle = \color{purple},
  sensitive=false
}
\lstdefinestyle{minileng}{
  language={minileng},
  basicstyle={\ttfamily\small},
  commentstyle={\small\itshape\color{Maroon}},
  ndkeywordstyle={\small},
  stringstyle={\ttfamily\color{red}},
  % frame={tb},
  breaklines=true,
  columns=[l]{fullflexible},
  % numbers=left,
  xrightmargin=0em,
  xleftmargin=3.5em,
  numberstyle={\scriptsize},
  stepnumber=1,
  numbersep=1em,
  lineskip=-0.5ex,
  keywordstyle = [1]\bfseries\color{blue},
  keywordstyle = [2]\bfseries\color{OliveGreen},
  keywordstyle = [3]\bfseries\color{cyan},
  keywordstyle = [4]\bfseries\color{orange},
  literate=
    {á}{{\'a}}1
    {Á}{{\'A}}1
    {à}{{\`a}}1
    {ã}{{\~a}}1
    {é}{{\'e}}1
    {É}{{\'E}}1
    {ê}{{\^e}}1
    {í}{{\'i}}1
    {ó}{{\'o}}1
    {Ó}{{\'O}}1
    {õ}{{\~o}}1
    {ú}{{\'u}}1
    {Ú}{{\'U}}1
    {ü}{{\"u}}1
    {ñ}{{\~n}}1
    {ç}{{\c{c}}}1,
    captionpos=b
}
\lstdefinestyle{text}{
  basicstyle={\ttfamily\small},
  breaklines=true,
  % columns=[l]{fullflexible},
  columns=flexible,
  % xrightmargin=0em,
  % xleftmargin=3.5em,
  lineskip=-0.5ex,
  literate=
    {á}{{\'a}}1
    {Á}{{\'A}}1
    {à}{{\`a}}1
    {ã}{{\~a}}1
    {é}{{\'e}}1
    {É}{{\'E}}1
    {ê}{{\^e}}1
    {í}{{\'i}}1
    {ó}{{\'o}}1
    {Ó}{{\'O}}1
    {õ}{{\~o}}1
    {ú}{{\'u}}1
    {Ú}{{\'U}}1
    {ü}{{\"u}}1
    {ñ}{{\~n}}1
    {ç}{{\c{c}}}1,
    captionpos=b
}
\lstdefinestyle{java}{
language=java,
tabsize = 4, %% set tab space width
showstringspaces = false, %% prevent space marking in strings, string is defined as the text that is generally printed directly to the console
% numbers = left, %% display line numbers on the left
commentstyle = \color{OliveGreen}, %% set comment color
keywordstyle = \color{blue}, %% set keyword color
stringstyle = \color{red}, %% set string color
rulecolor = \color{black}, %% set frame color to avoid being affected by text color
basicstyle = \small \ttfamily , %% set listing font and size
breaklines = true, %% enable line breaking
numberstyle = \tiny,
}

\renewcommand{\lstlistingname}{Código}

\lstnewenvironment{minileng}[1][]
    {\hfill\lstset{style=minileng,#1}\minipage{0.8\linewidth}}
    {\endminipage\hfill\null}

\lstnewenvironment{codigo}[1][]
    {\hfill\lstset{style=text,#1}\minipage{0.9\linewidth}}
    {\endminipage\hfill\null}


\usepackage{fancyvrb}

\usepackage{booktabs} % Allows the use of \toprule, \midrule and \bottomrule in tables for horizontal lines
\usepackage{graphicx}
\graphicspath{{imagenes/}{../imagenes/}} % Declara dos paths, uno con respecto a este documento y otro respecto a la carpeta secciones

\usepackage[margin=3cm, headheight=13.6pt]{geometry}
\usepackage{fancyhdr}
\fancyhf{}
\lhead{\emph{Procesadores de Lenguajes -- 2019-20}}
\rhead{\emph{Compilador de MiniLeng}}
\fancyfoot[C]{\thepage}

\usepackage[titletoc, title]{appendix}

\usepackage{fancyvrb}

\begin{document}

\begin{titlepage}
    \centering
    \includegraphics[width=1 cm]{logoUZ.jpg}

    \textsc{\large Universidad de Zaragoza}
    \rule{\textwidth}{1.6pt}\vspace*{-\baselineskip}\vspace*{2pt} % Thick horizontal rule
    \rule{\textwidth}{0.4pt} % Thin horizontal rule

    \vfill

    {\LARGE \scshape Procesadores de Lenguajes}

    \vspace{2cm}

    {\bfseries \Huge Compilador de MiniLeng}

    \vspace{.5cm}

    {\bfseries \Large Informe sobre el desarrollo del compilador}

    \vspace{3cm}



    % {\large \scshape Autor}

    %\vspace{0.2cm}

    {\large Fernando Peña Bes (NIA: 756012)}

    \vfill

    \large{Zaragoza, España}

    {Curso 2019\,--\,2020}

    \vfill

    \includegraphics[width=5.0cm]{EINA.png}

\end{titlepage}

\pagenumbering{Roman}


\vspace*{1cm}
\section*{\hfil Resumen \hfil}
\addcontentsline{toc}{section}{Resumen}

En este informe de documenta el desarrollo de un compilador para el lenguaje MiniLeng durante la asignatura Procesadores de Lenguajes.

MiniLeng es un lenguaje procedural, estructurado y fuertemente tipado, cuya sintaxis está inspirada en Pascal y C.

El lenguaje es \it{case-insensitive} y soporta comentarios tanto de una sola línea como multilínea. Se pueden declarar variables tanto globales como locales de tres tipos de dato: entero, carácter y booleano, y permite la declaración de acciones (procedimientos) anidadas con paso de parámetros tanto por valor como por referencia.

Se pueden utilizar las estructuras de control ``mientras que'' (while) y ``selección'' (if), y se soporta la asignación de valores a variables y la evaluación de expresiones aritméticas y lógicas.

El lenguaje tiene dos procedimientos predefinidos: ``escribir'', que permite escribir por pantalla variables simples y cadenas de caracteres constantes y ``leer'', que permite leer valores introducidos por el usuario y asignarlos a variables. También tiene dos funciones que permiten transformar un entero a carácter y viceversa: ``entacar'' y ``caraent''.

Por último, se ha añadido la posibilidad de trabajar con vectores unidimensionales. Se permite la declaración de vectores, la asignación de valores a las componentes, la asignación directa entre vectores del mismo tamaño, el uso de componentes en expresiones y el paso de componentes y vectores completos como argumentos al invocar una acción.

Un ejemplo de programa de MiniLeng válido semánticamente se puede ver en el Código~\ref{ejemplo}.

%%%%%% INSERTAR PROGRAMA
\begin{minileng}[caption={Ejemplo de programa en MiniLeng.},label={ejemplo}]
programa factorialVect;
    entero i, r, v[8];

    accion factorial(val entero n; ref entero r);
    principio
        si n < 2 ent
            r := 1;
        si_no
            factorial(n - 1, r);
            r := n * r;
        fsi
    fin

principio
    i := 0;
    mq i < 8
        factorial(i, v[i]);
        escribir("v[", i, "] = ", v[i], "\n");
        i := i + 1;
    fmq
fin\end{minileng}


Cada vez que se produce un error de compilación, el compilador muestra un mensaje explicativo junto al token más relevante en el error y su número de fila y columna en el programa. Distingue entre errores léxicos, sintácticos y semánticos. Cuando el compilador encuentra un error léxico termina la ejecución, pero cuando el error es sintáctico o semántico, intenta seguir compilando hasta el final del fichero para mostrar la máxima información posible sobre los errores que puede en el programa. También se incluyen avisos de compilación, como \textit{underflow}/\textit{overflow} y divisiones por cero.

Si no se han producido errores durante la compilación del programa (no se tienen en cuenta los avisos), se genera un fichero con el mismo nombre del fichero de entrada y extensión \tt{.code} con el código intermedio que contiene las instrucciones para ejecutar el programa en la máquina P. La generación de código se realiza utilizando el esquema AST, que da flexibilidad a la generación de código y permite introducir optimizaciones interesantes.

Para implementar el compilador, se ha utilizado JavaCC, un metacompilador de código abierto para el lenguaje Java.
Se ha desarrollado principalmente, sobre Eclipse 2020-06 con el plugin JavaCC Eclipse Plug-in 1.5.33.
\pagebreak


\tableofcontents
%\listoffigures
%\listoftables

\pagebreak

\setcounter{page}{1}
\pagenumbering{arabic}

\pagestyle{fancy}

\section{Análisis léxico}
El primer paso al construir el compilador fue implementar el analizador léxico. Su trabajo es tomar el programa como una cadena de caracteres y devolver un conjunto de \it{tokens} o símbolos, que son las unidades léxicas del lenguaje.

\subsection{Expresión regular utilizada para los identificadores}
% Detallar y explicar la expresión regular utilizada para los identificadores.
Según la especificación del lenguaje, los identificadores están compuestos por letras, números y el símbolo \car{\_}. No pueden comenzar por número y no pueden terminar en \car{\_}.

Una posible expresión regular que reconoce identificadores es la siguiente:

\begin{codigo}
< #DIGITO : [ "0"-"9" ] >
< #LETRA : [ "a"-"z" ] >
< tIDENTIFICADOR :
  (< LETRA > | "_"(< LETRA > | < DIGITO >))(("_")?(< LETRA > | < DIGITO >))* >
\end{codigo}

\subsection{Procesamiento de comentarios}
% Detallar y explicar cómo se han procesado los comentarios simples y multilínea.
En el lenguaje se permiten definir comentarios de una línea (simples) y multilínea.  Los comentarios de una línea se indican con el carácter \car{\%}, todo lo que haya desde ese carácter hasta el final de la línea se considerará como comentario. Los comentarios multilínea empiezan y terminan con la secuencia \car{\%\%}.

Para implementar los comentarios en JavaCC se han utilizado contextos léxicos.
Inicialmente el analizador se encuentra en el contexto \tt{DEFAULT}.
Se puede indicar el contexto al que pertenece una declaración léxica entre los signos \car{\textless} y \car{\textgreater}, de forma que sólo se tendrá en cuenta cuando el contexto actual sea ese. Para realizar un cambio de contexto después de reconocer una expresión utiliza el signo \car{:}.

El objetivo es filtrar los comentarios para que no sean enviados al analizador sintáctico. JavaCC define la declaración léxica \tt{SKIP} que permite justamente descartar la entrada reconocida. Teniendo estos aspectos en cuenta, la implementación propuesta para procesar comentarios es la siguiente:

\begin{codigo}
SKIP : {  " "| "\r"| "\t"| "\n"
| < "%"~["%"] > : COMENTARIO
| "%%"          : MULTICOMENTARIO}

// Ignorar comentarios en la entrada
<COMENTARIO> SKIP : {
  "\n" : DEFAULT
}

<COMENTARIO> MORE : {
   < ~[] >
}

<MULTICOMENTARIO> SKIP : {
  "%%" : DEFAULT
}

<MULTICOMENTARIO> MORE : {
   < ~[] >
}
\end{codigo}

Las expresiones reconocidas en las declaraciones \tt{MORE} son añadidas como prefijo a la siguiente expresión reconocida. En este caso se ha utilizado la expresión \verb|< ~[] >| dentro de declaraciones \tt{MORE} para descartar la entrada mientras se está dentro de un contexto de comentario.

\subsection{Implementación del modo \it{verbose}}
% Detallar y explicar cómo se ha implementado el modo verbose
El compilador dispone de una opción \it{verbose}. Cuando está activada, se muestra al final de la compilación una trabla con el número de ocurrencias de cada token que se han encontrado en el programa. La tabla tiene la siguiente forma:

\begin{codigo}[xleftmargin=3.5cm]
 +-------------------------------------+
 | Número de ocurrencias de los tokens |
 +-------------------------------------+
 | Token                         num.      |
 +-------------------------------------+
 | Palabras reservadas                     |
 |     PROGRAMA                     1      |
 |     PRINCIPIO                    3      |
 |     FIN                          3      |
 |     SI                           3      |
 |     ENT                          3      |
 |     FSI                          3      |
 | Caracteres de agrupación                |
 |     PARENTESIS_IZQ              27      |
 |     PARENTESIS_DER              27      |
 | Operadores aritméticos                  |
 |     SUMA                         2      |
 |     RESTA                        1      |
 | Operadores lógicos                      |
 |     AND                          2      |
 |     OR                           1      |
 | Tipos de dato                           |
 |     BOOLEANO                     1      |
 |     CARACTER                     4      |
 +-------------------------------------+
\end{codigo}

Para implementarla se ha creado una clase llamada \tt{TablaOcurrencias} en el paquete \url{lib.lexico}. Esta clase contiene un tipo enumerado para cada tipo de token (Palabras reservadas, Operadores lógicos, Tipos de dato, Valores...) con los nombres de los tokens que pertenecen al tipo. Cada tipo tiene asociado un vector con un contadores para cada uno de sus tokens, que se van incrementando conforme se reconocen en el código. Para incrementar los contadores, la clase define una función \tt{incrementar} que se llama desde el analizador léxico pasándole el token cuyo contador se quiere incrementar.

Por ejemplo, la expresión del analizador léxico que reconoce la palabra reservada \tt{PROGRAMA}, realiza la siguiente llamada a la clase:

\begin{codigo}
TOKEN : /* PALABRAS RESERVADAS */
{
  < tPROGRAMA : "programa" >
  {
    minilengcompiler.tabla_ocurrencias.incrementar(TablaOcurrencias.Reservadas.tPROGRAMA);
  }
  ...
}
\end{codigo}

\subsection{Gestión de errores léxicos}
% Detallar y explicar la función de error léxico implementada
Si mientras se está procesando un programa se encuentra un error léxico, el compilador termina y muestra un error como el siguiente:

\begin{codigo}
ERROR LÉXICO (<línea>, <columna>): símbolo no reconocido: `...' (\u...)
\end{codigo}

Si el carácter que ha producido el error no es ASCII, se muestra escapado con su código Unicode.

En JavaCC, los errores léxicos son capturados como excepciones de tipo \tt{Error} durante la ejecución del compilador. Cuando se captura esta excepción, se toma el carácter que ha producido el error junto con su fila y columna en la entrada y se muestran utilizando una función personalizada (definida en la clase \tt{ErrorLexico} de \url{lib.lexico}).

Esta función (Código~\ref{errorlexico}) toma como parámetros la línea y columna en la que ha ocurrido el error, el carácter erróneo y el carácter escapado a código Unicode, e imprime el error con el formato mencionado.

\begin{codigo}[style=java, keywordstyle=\color{blue}, stringstyle=\color{red},
    caption={Función \tt{ErrorLexico.deteccion}.},
    label={errorlexico}]

public static void deteccion(int linea, int columna, String error, String error_escapado) {
    contadorErrores++;

    // Imprime el caracter erroneo, si puede no ser imprimible por la terminal,
    // imprime su equivalente escapado.
    System.err.println("ERROR LÉXICO (" + linea +
        ", " + (columna - 1) + "): " +
        "Caracter no reconocido: '" + error + "'" +
        (!error.equals(error_escapado) ? " (" + error_escapado + ")" : "")
    );
}
\end{codigo}

\documentclass[../main.tex]{subfiles}

\begin{document}
\section{Análisis sintáctico}

\subsection{Modificaciones sobre la gramática propuesta y decisiones de diseño}
% Modificaciones sobre la gramática propuesta y decisiones de diseño adoptadas.
En la gramática propuesta, algunas reglas estaban incompletas.
A continuación, se detalla como se han completado y cuáles han sido las decisiones de diseño adoptadas.

\begin{description}
\item[Parametros formales] Contiene la definición de los parámetros de una acción, esta estructura es opcional por lo que se ha utilizado el símbolo \car{?}:
\begin{verbatim}
parametros_formales	::=	( parentesis_izq ( lista_parametros )? parentesis_der )?

lista_parametros	::=	parametros ( fin_sentencia parametros )*

\end{verbatim}


\item[Lista de sentencias]

\begin{verbatim}

lista_sentencias	::=	sentencia ( sentencia )*

sentencia	::=	( leer | escribir | identificacion | seleccion | mientras_que )


leer	::=	<tLEER> parentesis_izq lista_asignables parentesis_der fin_sentencia

lista_asignables	::=	identificadores

escribir	::=	<tESCRIBIR> parentesis_izq lista_escribibles parentesis_der fin_sentencia
lista_escribibles	::=	lista_expresiones




seleccion	::=	<tSI> expresion <tENT> lista_sentencias ( <tSI_NO> lista_sentencias )* <tFSI>

\end{verbatim}


\item[Expresiones]

\begin{verbatim}

lista_expresiones	::=	expresion ( sep_variable expresion )*

expresion	::=	expresion_simple ( operador_relacional expresion_simple )?
operador_relacional	::=	( <tIGUAL> | <tMENOR> | <tMAYOR> | <tMAI> | <tMEI> | <tNI> )
expresion_simple	::=	( <tMAS> | <tMENOS> )? termino ( operador_aditivo termino )*
operador_aditivo	::=	( <tMAS> | <tMENOS> | <tOR> )
termino	::=	factor ( operador_multiplicativo factor )*
operador_multiplicativo	::=	( <tPRODUCTO> | <tDIVISION> | <tMOD> | <tAND> )

\end{verbatim}

\end{description}

Los símbolos \car{,}, \car{;}, \car{(} y \car{)}  se han sustituido en la gramática por sus tokens correspondientes: \tok{tSEP\_VARIABLE}, \tok{tFIN\_SENTENCIA}, \tok{tPRENTESIS\_IZQ} y \tok{tPARENTESIS\_DER}, de esta forma sería posible cambiar los caracteres en un futuro de forma sencilla.

\subsection{Gestión de errores y avisos (\it{warnings}) sintácticos}
% Detallar y explicar el manejo de excepciones y la función de error sintáctico implementada.

\end{document}

\section{Tabla de símbolos}
\label{tabla}

La tabla de símbolos es una estructura auxiliar que se utiliza principalmente durante las fases de análisis semántico y generación de código. En ella se almacena la información asociada a los símbolos declarados en el programa. De esta forma, cuando se utiliza un símbolo durante el programa, es posible conocer si está definido, cuál es su tipo, etc.

Para implementarla se ha utilizado una tabla hash en la que las colisiones se resuelven mediante listas enlazadas.

\subsection{Simbolos}
El primer paso fue definir lo que es un símbolo. Para ello se creó la clase \url{lib.semantico.Simbolo} cuyos atributos se pueden ver en el Código~\ref{simbolo}.

\begin{codigo}[style=java,caption={Atributos de la clase \url{lib.semantico.Simbolo}.},label={simbolo}]
public enum Tipo_simbolo {
    PROGRAMA, VARIABLE, ACCION, PARAMETRO
};

// Representa el tipo de variable
public enum Tipo_variable {
    DESCONOCIDO, ENTERO, BOOLEANO, CHAR, CADENA
};

// Representa la clase de los parámetros en las acciones
public enum Clase_parametro {
    VAL, REF
};

/////////////////////////////
// Atributos del símbolo //
/////////////////////////////

String nombre;
Integer nivel; // Nivel en el que se ha declarado el símbolo (primer nivel = 0)
Integer dir; // Dirección del símbolo

Tipo_simbolo tipo;
Tipo_variable variable;
Clase_parametro parametro;

ArrayList<Simbolo> lista_parametros; // Lista de símbolos que representan los parámetros de una acción

Boolean vector = false; // Vale true si el símbolo es una variable o parametro vector
Integer longitud; // Longitud para los vectores

Boolean inicializado = false; // Vale true si el símbolo es una variable o parámetro y ha sido inicializado
\end{codigo}

Para representar los vectores no se han definido tipos nuevos de símbolo, se ha decidido añadir un booleano (\tt{vector}) que valdrá \tt{true} si el símbolo es una variable o parámetro vector y un entero, que almacenará la longitud del vector.

\subsection{Función de hash}
% Detallar y explicar la función de hash implementada
Como función de hash se ha escogido la función de Pearson. Dada una cadena, es capaz de calcular un hash rápidamente, no hay clases de entradas más propensas a colisiones que otras, las cadenas que difieren en un carácter nunca colisionan y tampoco se generan colisiones entre anagramas.

En su versión más sencilla esta función genera hashes de 8 bits (0-255), y requiere un array de enteros con los números del 0 al 255 ordenados aleatoriamente.

No se espera que el compilador trabaje con programas muy complejos y con gran cantidad identificadores diferentes, así que esta función es adecuada. Su implementación se puede ver en el Código~\ref{hash}.

\begin{codigo}[style=java,caption={Función \url{lib.semantico.TablaSimbolos.h}.},label={hash}]
/*
 * Función de hash utilizando el algoritmo de Pearson
 */
private int h(String cadena) {
    int h = 0;
    for (int i = 0; i < cadena.length(); i++) {
        h = T[h ^ cadena.charAt(i)];
    }
    return h;
}
\end{codigo}

Para generar el vector \tt{T}, se incializa un vector de 256 componentes con los números de 0 al 255 en orden y se permutan utilizando el algoritmo Fisher-Yates. Este algoritmo se muestra en el código~\ref{mezcla}.

\begin{codigo}[style=java,caption={Función \url{lib.semantico.TablaSimbolos.mezclaVector}.},label={mezcla}]
/*
 * Algoritmo de Fisher-Yates para permutar aleatoriamente los elementos de un
 * vector
 */
private void mezclaVector(int[] a) {
    Random rnd = new Random();
    for (int i = a.length - 1; i > 0; i--) {
        // Genera un número aleatorio j tal que 0 <= j 0 <= i
        int j = rnd.nextInt(i + 1);

        // Intercambia a[j] y a[i]
        int aux = a[j];
        a[j] = a[i];
        a[i] = aux;
    }
}
\end{codigo}


\subsection{Implementación de la tabla y manejo de colisiones}
% Detallar y explicar la implementación de la tabla y el manejo de colisiones.
%
La tabla hash se ha implementado como un array de listas enlazadas. Los nuevos elementos se concatenan al principio de las listas correspondientes a sus hashes. Por tanto, cuando se produce una colisión, el elemento más recientemente introducido será el que primero esté en la lista.

En MiniLeng es posible declarar identificadores con el mismo nombre en niveles distintos, de forma que las declaraciones locales oculten a las globales. Para encontrar cuál es el último símbolo declarado con un cierto nombre, basta con calcular su hash y recorrer desde el principio la lista enlazada correspondiente al hash hasta encontrar el primer símbolo cuyo nombre es el que se estaba buscando. Cada vez que un símbolo es ocultado por otro con el mismo nombre, se muestra un aviso para que el programador lo tenga en cuenta.

En cada nivel, hay un único espacio de nombres. No se permite introducir símbolos con el mismo nombre en el mismo nivel, ya que habría una ambigüedad a la hora de decidir a qué símbolo se está haciendo referencia en el código.

Se pensó en implementar un espacio de nombres distinto para las acciones y las variables/parámetros, y distinguir el tipo de símbolo según el uso que se le de en una sentencia. Sin embargo, esta alternativa puede llevar a varios problemas. Por ejemplo, si el lenguaje soportara funciones y la sintaxis para definir el valor a devolver fuera asignarlo al símbolo de la función, no se podría saber si realmente se está asignando un valor a la función o a una variable con el mismo nombre. Este problema se podría solucionar dando al programador la opción de crear y especificar espacios de nombres como en Pascal o C, aunque esto aumentaría la complejidad de la sintaxis del lenguaje.

En los Códigos~\ref{introducir} y~\ref{buscar} se muestra la implementación de las funciones para insertar y buscar símbolos en la tabla de símbolos, siguiendo las ideas anteriores.

\begin{codigo}[style=java,caption={Función \url{lib.semantico.TablaSimbolos.introducir_simbolo}},label={introducir}]
/*
 * Si existe un símbolo en la tabla del mismo nivel y con el mismo, nombre,
 * lanza una excepción. De lo contrario, introduce el símbolo pasado como
 * parámetro.
 */
private Simbolo introducir_simbolo(Simbolo simbolo) throws SimboloYaDeclaradoException {
    int clave = h(simbolo.getNombre());

    for (Simbolo s : tabla_hash[clave]) {
        // Si el símbolo ya está declarado en ese mismo nivel, lanzar una excepción
        if (s.getNombre().equals(simbolo.getNombre()) && s.getNivel() == simbolo.getNivel()) {
            throw new SimboloYaDeclaradoException();
        }
        // Si hay un símbolo ya declarado con ese nombre en otro nivel, mostrar un aviso
        else if (s.getNombre().equals(simbolo.getNombre())) {
            Token t = minilengcompiler.token;
            if (simbolo.ES_VECTOR()) {
                // Evitar que el token sea ']'
                t.image = simbolo.nombre;
            }
            Aviso.deteccion("Este símbolo, definido en el nivel " + simbolo.nivel +
                    ", va a ocultar a otro definido con el mismo nombre en el nivel " + s.nivel + "",
                    t);
        }
    }

    // Si no, se añade
    tabla_hash[clave].addFirst(simbolo);

    return simbolo;
}
\end{codigo}


\begin{codigo}[style=java,caption={Función \url{lib.semantico.TablaSimbolos.buscar_simbolo}},label={buscar}]
/*
 * Busca en la tabla el símbolo de mayor nivel cuyo nombre coincida con el del
 * parámetro (se distinguen minúsculas y mayúsculas). Si existe, devuelve un
 * puntero como resultado, de lo contrario lanza una excepción.
 */
public Simbolo buscar_simbolo(String nombre) throws SimboloNoEncontradoException {
    int clave = h(nombre);
    for (Simbolo s : tabla_hash[clave]) {
        if (s.nombre.equals(nombre)) {
            return s;
        }
    }

    // Si no se ha encontrado
    throw new SimboloNoEncontradoException();
}
\end{codigo}

\documentclass[../main.tex]{subfiles}

\begin{document}
\section{Análisis semántico}
\subsection{Gestión de errores y avisos (\it{warnings}) semánticos}
% Detallarlos errores y avisos (warnings) que se han implementado

\subsection{Controles sobre los tipos de datos}
% Detallar y explicar los controles que se han implementado sobre los tipos de datos (desbordamientos, booleanos, etc.)

\subsection{Propagación y manejo de valores constantes en las expresiones}
% Detallar y explicar cómo se ha implementado la propagación y el manejo de valores constantes en las expresiones.

\subsection{Comprobaciones sobre parámetros \tt{VAL} y \tt{REF}}
% Detallar y explicar cómo se han implementado las comprobaciones sobre parámetros VAL/REF. ¿Dónde se realizan estas comprobaciones? Adjuntar ejemplos que validen la funcionalidad implementada.


No se van a usar parámetros ocultos. Como estamos trabajando con objetos Java, las acciones pueden
tener una lista de referencias a los parámetros. Los parámetros se borrarán de la tabla en el
momento que la función deje de ser visible y se borre.

\end{document}

\section{Generación de código}
% Detallar y justificar la elección del esquema de generación de código utilizado (secuencial/AST).
Esta es la última fase del compilador. Tras su ejecución se genera un fichero con las instrucciones para ejecutar el programa de MiniLeng en la máquina P.

\subsection{Justificación del esquema escogido}\label{eleccion}
Para la fase de generación de código se ha escogido el esquema AST. En esta implementación no se genera el árbol abstracto de sintaxis de forma explicita, si no que se aprovecha el árbol de análisis generado por el analizador sintáctico para mejorar el código final.

Este esquema es mucho más flexible que el secuencial, ya que permite reordenar las instrucciones del programa y eliminar las que no sean necesarias. Cada instrucción del programa para la máquina P se almacena como una cadena de caracteres y estás sólo se imprimen al final de la compilación si no se han encontrado errores en el código.

Las cadenas que representan las instrucciones se van guardando en listas. Se ha creado una clase llamada \tt{ListaInstr} en el paquete \url{lib.generacioncodigo} con la que se pueden gestionar estas listas. El único atributo de la clase es un \tt{ArrayList<String>} y se incluyen métodos para añadir en la lista todas las instrucciones necesarias para generar el código final.

Los métodos introducen bloques de instrucciones que realizan una determinada función. Se llaman desde el lugar correspondiente de la gramática para generar el conjunto de instrucciones que se necesite.

La idea es que las reglas de la gramática manipulan y devuelven las listas de instrucciones hasta que todas llegan a la regla inicial (\tt{programa}). En este lugar se juntan las listas devueltas y se genera una lista final con todas las instrucciones del programa.

Esta es una interfaz limpia y fácil de modificar, ya que casi toda la generación de código esta centralizada en una clase. Por ejemplo, el siguiente método (Código~\ref{addseleccion}) genera el código necesario para añadir una selección a la lista.

\begin{codigo}[style=java, caption={Método \url{lib.generacioncodigo.ListaInstr.addSeleccion}}, label={addseleccion}]
public void addSeleccion(ListaInstr expresion, ListaInstr si, ListaInstr sino, Integer etiq1, Integer etiq2) {
    // El código generado tiene la siguiente forma:
    //
    //			<expresion>
    //			JMF ETIQ1
    //			<si>
    //			JMP ETIQ2
    // ETIQ1:
    //			<sino>
    // ETIQ2:
    addComentario("Seleccion");
    concatenarLista(expresion);
    addSaltoFalse(etiq1);
    addComentario("Si");
    concatenarLista(si);
    addSaltoIncod(etiq2);
    addEtiqueta(etiq1);
    addComentario("Si no");
    concatenarLista(sino);
    addEtiqueta(etiq2);
}
\end{codigo}

Toma la listas de instrucciones necesaria para calcular la condición de la selección, además del código que hay que ejecutar en el `si' y en el `si\_no'.

Si en el análisis semántico se detecta que no hay bloque `si\_no', se puede generar el código utilizando el método del Código~\ref{addseleccionsimple} para evitar el salto incondicional.

\begin{codigo}[style=java, caption={Método \url{lib.generacioncodigo.ListaInstr.addSeleccionSimple}}, label={addseleccionsimple}]
public void addSeleccionSimple(ListaInstr expresion, ListaInstr si, Integer etiq) {
    // El código generado tiene la siguiente forma:
    //
    //			<expresion>
    //			JMF ETIQ
    //			<si>
    // ETIQ:
    addComentario("Seleccion simple");
    concatenarLista(expresion);
    addSaltoFalse(etiq);
    addComentario("Si");
    concatenarLista(si);
    addEtiqueta(etiq);
}
\end{codigo}

Esta flexibilidad para mejorar el código dependiendo del programa de entrada no se tiene al utilizar un esquema de generación de código secuencial.

El código que se encarga de gestionar el nivel actual durante la compilación, los números de las etiquetas y las direcciones del bloque de activación, además de la escritura de la lista final en el fichero de salida, se encuentra en la clase \tt{GeneracionCodigo}, también en el paquete \url{lib.generacioncodigo}.

\documentclass[../main.tex]{subfiles}

\begin{document}
\section{Mejoras introducidas}
% Detallar las mejoras y puntos opcionales que se han desarrollado sobre la propuesta inicial.
% Especificar detalles de implementación, decisiones de diseño y pruebas/ejemplos de las funcionalidades añadidas.
% Separar las mejoras en bloques, según correspondan al análisis léxico, semántico o a la generación de código.

\subsection{Mejoras en el comportamiento del compilador}

Errores en el uso del programa

\begin{itemize}
\item Si se utiliza una opción inválida o no se especifican correctamente, se muestra el siguiente error:
\begin{codigo}{text}
MiniLeng: Opción inválida <error>
\end{codigo}

\item Si el fichero a compilar no tiene extensión \tt{.ml}, se muestra el error:
\begin{codigo}{text}
MiniLeng: El fichero a compilar tiene que tener extensión .ml
          Fichero introducido: '...'
\end{codigo}

\item Si no se encuentra el fichero a compilar, se indica de la siguiente forma:
\begin{codigo}{text}
MiniLeng: No se ha encontrado el fichero '...'
\end{codigo}

\end{itemize}

Al final de la compilación se muestra un recuento de los errores encontrados:

\begin{codigo}{text}
Errores léxicos: ...
Errores sintácticos: ...
Veces activado panic mode: ...
\end{codigo}

Y se muestra

\begin{codigo}{text}
No se ha podido compilar el programa
\end{codigo}

si ha habido errores durante la compilación, o:

\begin{codigo}{text}
Compilado sin errores
\end{codigo}

si la compilación ha sido exitosa.

Si durante la compilación se ha activado el modo pánico, se muestra el siguiente mensaje:

\begin{codigo}{text}
Se ha activado el panic mode durante la compilación. Corrige los errores y vuelve a compilar.
\end{codigo}

para indicar al usuario que tiene que hace falta volver a compilar el programa.



\subsection{Mejoras en el análisis léxico}
La opción tokens muestra por pantalla los nombres de los tokens que se van reconociendo conforme se analiza el programa. Ejemplo:

\begin{codigo}{text}
tPROGRAMA
tIDENTIFICADOR (Valor: mi_programa)
tFIN_SENTENCIA
tENTERO
tIDENTIFICADOR (Valor: n)
tFIN_SENTENCIA
...
tFIN
\end{codigo}

\subsection{Mejoras en el análisis semántico}
  - Se permiten acciones anidadas con cualquier nivel de profuncidad. Las declaraciones de acciones
    anidadas deben entre la zona de declaración de las variables y la palabra 'principio' de la
    función padre:

  \begin{verbatim}
      accion accion1;
        % Declaración de variables de accion1
        accion accion2;
          % Declaración de variables de accion2
          accion accion3;
            % Declaración de variables de accion3
          principio
            % Sentencias de accion3
          fin
        principio
          % Sentencias de accion2
        fin
      principio
        % Sentencias de accion1
      fin
  \end{verbatim}

  - Se permiten bloques 'seleccion' y 'mientras que' anidados con cualquier nivel de profundidad:

  \begin{verbatim}
      SI <condicion> ENT
        % Sentencias
        SI <condicion> ENT
          % Sentencias
        SI_NO
          % Sentencias
        FSI
        % Sentencias
      SI_NO
        % Sentencias
      FSI

      MQ <condicion>
        % Sentencias
        MQ <condicion>
          % Sentencias
        FMQ
        % Sentencias
      FMQ
  \end{verbatim}

  - Cualquier bloque ('principio/fin', 'seleccion' o 'mientras que') debe contener al menos una
    sentencia.




\subsubsection{Panic Mode}
El compilador dispone de la opción panic (\tt{-p}), que activa el modo pánico durante la compilación. Se entra en este modo cada vez que se produce un error sintáctico porque se esperaba un punto y coma. El analizador descarta todos los tokens siguientes al error hasta que encuentra un carácter \car{;} o el fichero se acaba, entonces sale del modo y sigue analizando la entrada.

Cuando se entra en el modo pánico, se informa al usuario con las siguientes líneas:

\begin{codigo}{text}
MiniLeng: ERROR SINTÁCTICO (línea ..., columna ...): Token incorrecto: '...'. Se esperaba ';'
PANIC MODE: Iniciado panic mode
\end{codigo}

Y se va mostrando en una línea diferente cada token descartado:

\begin{codigo}{text}
> PANIC MODE: Token descartado: '...'
\end{codigo}

Cuando se sale del modo, se muestra:

\begin{codigo}{text}
  > PANIC MODE (línea ..., columna ...): Se ha encontrado ';'
PANIC MODE: Terminado panic mode
\end{codigo}

\subsection{Mejoras en la generación de código}

\end{document}

\section{Pruebas realizadas}
En esta sección se detalla la metodología seguida para verificar el funcionamiento del compilador desarrollado en cada fase. Para cada fase se incluyen algunos tests junto a sus resultados, no se tratan de unas pruebas intensivas pero sirven para dar una idea del proceso que se ha seguido en la fase de pruebas.

\subsection{Analizador léxico}
El objetivo del analizador léxico es que reconozca todos los tokens de la entrada y que se lancen errores léxicos cuando se encuentran caracteres no soportados. Para probarlo, se intentó compilar varios programas diferentes activando la opción \it{tokens} (\tt{-t}, \tt{-{}-tokens}), que muestra los tokens que se van reconociendo en en análisis, para verificar que todos son reconocidos correctamente.

\subsubsection{Test 1}
% Poner un programa un poco complicado y ver que todos los tokens salen bien reconocidos
Vamos a considerar que tenemos el siguiente programa:

\begin{codigo}[style=minileng,numbers=left]
programa factorial;
entero n, solucion;

accion fact (val entero numero; ref entero solucion);
entero solRecur;
principio
	si numero <= 1 ent
		solucion := 1;
	si_no
	  	fact(numero - 1, solRecur);
		solucion := solucion * solRecur;
	fsi
fin

principio
	n := 10;
	fact(n, solucion);
	escribir(n, "! es: ", solucion);
fin
\end{codigo}

Al compilar el programa utilizando la opción \tt{-t}, se obtiene la siguiente salida, que si se observa detenidamente contiene todos los tokens del programa:

\noindent\begin{minipage}{.45\textwidth}
\begin{codigo}
tPROGRAMA
tIDENTIFICADOR (Valor: factorial)
tFIN_SENTENCIA
tENTERO
tIDENTIFICADOR (Valor: n)
tSEP_VARIABLE
tIDENTIFICADOR (Valor: solucion)
tFIN_SENTENCIA
tACCION
tIDENTIFICADOR (Valor: fact)
tPARENTESIS_IZQ
tVAL
tENTERO
tIDENTIFICADOR (Valor: numero)
tFIN_SENTENCIA
tREF
tENTERO
tIDENTIFICADOR (Valor: solucion)
tPARENTESIS_DER
tFIN_SENTENCIA
tENTERO
tIDENTIFICADOR (Valor: solRecur)
tFIN_SENTENCIA
tPRINCIPIO
tSI
tIDENTIFICADOR (Valor: numero)
tMEI
tCONSTENTERA (Valor: 1)
tENT
tIDENTIFICADOR (Valor: solucion)
tOPAS
tCONSTENTERA (Valor: 1)
tFIN_SENTENCIA
tSI_NO
tIDENTIFICADOR (Valor: fact)
tPARENTESIS_IZQ
tIDENTIFICADOR (Valor: numero)
\end{codigo}
\end{minipage}\hfill
\begin{minipage}{.45\textwidth}
\begin{codigo}
tMENOS
tCONSTENTERA (Valor: 1)
tSEP_VARIABLE
tIDENTIFICADOR (Valor: solRecur)
tPARENTESIS_DER
tFIN_SENTENCIA
tIDENTIFICADOR (Valor: solucion)
tOPAS
tIDENTIFICADOR (Valor: solucion)
tPRODUCTO
tIDENTIFICADOR (Valor: solRecur)
tFIN_SENTENCIA
tFSI
tFIN
tPRINCIPIO
tIDENTIFICADOR (Valor: n)
tOPAS
tCONSTENTERA (Valor: 10)
tFIN_SENTENCIA
tIDENTIFICADOR (Valor: fact)
tPARENTESIS_IZQ
tIDENTIFICADOR (Valor: n)
tSEP_VARIABLE
tIDENTIFICADOR (Valor: solucion)
tPARENTESIS_DER
tFIN_SENTENCIA
tESCRIBIR
tPARENTESIS_IZQ
tIDENTIFICADOR (Valor: n)
tSEP_VARIABLE
tCONSTCAD (Valor: "! es: ")
tSEP_VARIABLE
tIDENTIFICADOR (Valor: solucion)
tPARENTESIS_DER
tFIN_SENTENCIA
tFIN
\end{codigo}
\end{minipage}

\begin{codigo}
\end{codigo}

\subsubsection{Test 2}
% Comprobar que se producen errores
Cuando se encuentra algún carácter no soportado por el lenguaje se produce un error léxico, a no ser que este se encuentre en un comentario, sea una constante carácter o forme parte de una cadena constante.

Si se renombra a \tt{solución} la variable \tt{solución} modificando la linea 2 por la siguiente:

\begin{codigo}[style=minileng]
entero n, solución
\end{codigo}

aparece el error léxico:

\begin{codigo}
ERROR LÉXICO (2, 16): Caracter no reconocido: `ó' (\u00f3)
\end{codigo}

\subsection{Analizador sintáctico}
Para probar el analizador sintáctico una buena estrategia es intentar compilar programas léxicamente correctos pero con errores sintácticos y comprobar que se el compilador lanza errores en los lugares adecuados.

\subsubsection{Test 1}
Tenemos el siguiente sencillo programa, al que le falta un punto y coma en la línea 4.

\begin{codigo}[style=minileng,numbers=left]
programa prueba_sintactico;
entero n;
principio
    n := 1 + 1
    escribir("Hola");
fin
\end{codigo}

Cuando se intenta compilar, el analizador léxico pasa al analizador sintáctico el token \tt{escribir} cuando se necesitaba \car{;}, por lo que aparece el siguiente error:

\begin{codigo}
ERROR SINTÁCTICO (7, 1): Token incorrecto: `escribir'. Se esperaba `;'

Errores sintácticos: 1

No se ha podido compilar el programa.
\end{codigo}

\subsubsection{Test 2}
En la siguiente declaración de acción no se han especificado la clase del primer parámetro:

\begin{codigo}[style=minileng,numbers=left,firstnumber=4]
accion miAccion(entero n, val caracter c);
principio
    escribir(n, c);
fin
\end{codigo}

Al compilar, se ve como el analizador sintáctico intenta seguir compilando el programa después de detectar el primer error, aunque esto lleve a más errores sintácticos:

\begin{codigo}
ERROR SINTÁCTICO (4, 17): Token incorrecto: `entero'. Se esperaba `('
ERROR SINTÁCTICO (4, 17): Token incorrecto: `entero'. Se esperaba `;'
ERROR SINTÁCTICO (4, 27): Token incorrecto: `val'. Se esperaba un identificador
ERROR SINTÁCTICO (4, 27): Token incorrecto: `val'. Se esperaba `;'
ERROR SINTÁCTICO (4, 27): Token incorrecto: `val'. Se esperaba el delimitador de principio de bloque: `principio'
ERROR SINTÁCTICO (4, 27): Token incorrecto: `val'. Se esperaba una sentencia
ERROR SINTÁCTICO (4, 27): Token incorrecto: `val'. Se esperaba el delimitador fin de bloque: `fin'
ERROR SINTÁCTICO (4, 27): Token incorrecto: `val'. Se esperaba el delimitador de principio de bloque: `principio'
ERROR SINTÁCTICO (4, 27): Token incorrecto: `val'. Se esperaba una sentencia
ERROR SINTÁCTICO (4, 27): Token incorrecto: `val'. Se esperaba el delimitador fin de bloque: `fin'
ERROR SINTÁCTICO (4, 27): Token incorrecto: `val'. La declaración del programa es incorrecta

Errores sintácticos: 11

No se ha podido compilar el programa.
\end{codigo}

\subsection{Analizador semántico}
Para probar el analizador semántico se siguió una estrategia similar a la del sintáctico. Se intentan compilar programas sintácticamente correctos pero con errores semánticos y se verifica que los errores aparezcan en los lugares esperados.

\subsubsection{Test 1}
El siguiente programa intenta declarar dos variables con el mismo nombre:

\begin{codigo}[style=minileng,numbers=left]
programa prueba_semantico;
entero n, n;

    accion miAccion
    entero n;
    principio
        escribir(n);
    fin

principio
    n := 1;
    escribir(n);
fin
\end{codigo}

Al intentar compilarlo se obtiene lo siguiente:

\begin{codigo}
ERROR SEMÁNTICO (2, 11): Símbolo: `n'. No se puede redefinir el símbolo
AVISO (5, 8): Símbolo: `n'. Este símbolo, definido en el nivel 1, va a ocultar a otro definido con el mismo nombre en el nivel 0

Errores semánticos: 1
Avisos: 1

No se ha podido compilar el programa.
\end{codigo}

El compilador detecta que se ha intentado declarar una variable dos veces en la línea 2. Dentro de la acción también se declara una variable llamada \tt{n}, pero al estar en un nivel inferior no se está redefiniendo la variable anterior, si no ocultándola.

\subsubsection{Test 2}
En el siguiente programa se intenta operar con tipos diferentes:

\begin{codigo}[style=minileng,numbers=left]
programa prueba_semantico;
entero n;
booleano b;
principio
    b := true or false;
    n := "a" + b;
fin
\end{codigo}

por lo que el compilador genera:

\begin{codigo}
ERROR SEMÁNTICO (6, 14): Símbolo: `+'. El operando 1 debe ser entero

Errores semánticos: 1

No se ha podido compilar el programa.
\end{codigo}

El término que contiene la suma devuelve un \tt{RegistroExpr} de tipo \tt{ENTERO}, por lo que el error no se propaga hasta la asignación.


\subsection{Generación de código}
Para probar la generación de código, se compilaron varios programas semánticamente correctos y se analizó el código generado. Este código también se probó a ejecutar en Hendrix utilizando el ensamblador e intérprete proporcionados. En esta fase ha sido de gran utilidad generar comentarios con la función de cada grupo de instrucciones, para saber que parte del código intermedio corresponde con cada parte del programa MiniLeng.

\subsubsection{Test 1}
El siguiente programa contiene una sentencia mientras que:

\begin{codigo}[style=minileng,numbers=left]
programa generacion_codigo;
entero i;
principio
    i := 10;
    mq i < 10
        escribir(i);
        i := i - 1;
    fmq
fin
\end{codigo}

Analizando el código generado se ve que es correcto y leyendo los comentarios se puede ver la correspondencia con cada parte del programa inicial.

\begin{codigo}
; Programa: generacion_codigo
	ENP L2
; Comienzo de programa: generacion_codigo
L2:
; Asignación a i
; Obtener dirección de VARIABLE ENTERO:          i [0, 3]
	SRF 0 3
; Sustituida expresion constante
	STC 10
	ASG
; Mientras que
	JMP L0
L1:
; Escribir entero
; Obtener valor de VARIABLE ENTERO:          i [0, 3]
	SRF 0 3
	DRF
	WRT 1
; Asignación a i
; Obtener dirección de VARIABLE ENTERO:          i [0, 3]
	SRF 0 3
; Obtener valor de VARIABLE ENTERO:          i [0, 3]
	SRF 0 3
	DRF
; Sustituida expresion constante
	STC 1
	SBT
	ASG
L0:
; Condicion mientras que
; Obtener valor de VARIABLE ENTERO:          i [0, 3]
	SRF 0 3
	DRF
; Sustituida expresion constante
	STC 10
	LT
	JMT L1
; Fin de generacion_codigo
	LVP
\end{codigo}

\subsubsection{Test 2}
El siguiente programa contiene varias expresiones constantes:

\begin{codigo}[style=minileng,numbers=left]
programa generacion_codigo;
entero e;
booleano b;
principio
    e := (20 + 4) * 5 / 4;
    b := (2 * 100) = (100 * 2);
fin
\end{codigo}

El compilador es capaz de calcularlas correctamente y almacenar sus valores directamente sobre la pila:

\begin{codigo}
; Programa: generacion_codigo
	ENP L0
; Comienzo de programa: generacion_codigo
L0:
; Asignación a e
; Obtener dirección de VARIABLE ENTERO:          e [0, 3]
	SRF 0 3
; Sustituida expresion constante
	STC 30
	ASG
; Asignación a b
; Obtener dirección de VARIABLE BOOLEANO:        b [0, 4]
	SRF 0 4
; Sustituida expresion constante
	STC 1
	ASG
; Fin de generacion_codigo
	LVP
\end{codigo}

\subsubsection{Test 3}
% Poner el código generado al acceder a una componente de vector a través de una variable en un parámetro por referencia
A continuación, se va a intentar acceder a una componente de vector que se ha pasado por referencia a una acción para comprobar que el código generado en este caso es correcto.

\begin{codigo}[style=minileng,numbers=left]
programa generacion_codigo;
entero v[10];

accion miAccion(ref entero v[10]);
principio
    escribir(v[1]);
fin

principio
    v[1] := 100;
    miAccion(v);
fin
\end{codigo}

Se ha generado el siguiente código. Se puede comprobar que es correcto y que sigue los aspectos de la generación de código explicados en las secciones anteriores.

\begin{codigo}
; Programa: generacion_codigo
	ENP L2
; Acción: miAccion
L0:
; Recuperar parametro 1: PARAMETRO REF ENTERO:     v[10] [1, 3]
	SRF 0 3
	ASGI
; Comienzo de acción: miAccion
; Escribir entero
; Obtener valor de PARAMETRO REF ENTERO:     v[10] [1, 3]
	SRF 0 3
	DRF
; Indice
; Sustituida expresion constante
	STC 1
	PLUS
	DRF
	WRT 1
	CSF
; Comienzo de programa: generacion_codigo
L2:
; Asignación a v
; Obtener dirección de VARIABLE ENTERO:          v[10] [0, 3]
; Indice: 1
	SRF 0 4
; Sustituida expresion constante
	STC 100
	ASG
; Apilar argumento 1: PARAMETRO REF ENTERO:     v[10] [1, 3]
; Obtener dirección de VARIABLE ENTERO:          v[10] [0, 3]
	SRF 0 3
; Invocar a ACCION:                   miAccion(v) [0, 0]
	OSF 14 0 L0
; Fin de generacion_codigo
	LVP
\end{codigo}


\renewcommand\appendixname{Anexo}

\begin{appendices}
\documentclass[../main.tex]{subfiles}

\begin{document}
\section{Gramática de MiniLeng implementada} \label{gramatica}


\end{document}
% \documentclass[../main.tex]{subfiles}

\begin{document}
\section{Uso del compilador}

\begin{verbatim}
Compilador de MiniLeng -- v2.2 (abril de 2020)
Autor: Fernando Peña Bes (NIA: 756012)

Uso: minilengcompiler [opciones] [fichero]

Opciones:
  -v, --verbose  Mostrar un resumen de los símbolos utilizados en el programa
  -p, --panic	 Compila con panic mode
  -t, --tokens   Muestra los tokens que se van reconociendo
  -h, --help	 Imprimir ayuda (esta pantalla) y salir
  --version      Imprimir información de la versión y salir


Ejemplos de uso:
    minilengcompiler
    minilengcompiler -v
    minilengcompiler -v fichero.ml
    minlengcompiler -p -v fichero.ml

El fichero con el programa a compilar tiene que tener extensión .ml
\end{verbatim}

\end{document}

\end{appendices}

\clearpage{}
\newpage
\nocite{*}
\addcontentsline{toc}{section}{Referencias}
\printbibliography

\end{document}
